% Outline: Methodology (1 page cap, >=2 tables; recipes + provenance, no results)
\section{Methodology}

\subsection{Data and Labels}
Training and offline evaluation use a 6{,}000-request stratified
sample of ToolACE tool-calling conversations, labeled with completion
lengths measured against Qwen3.5-9B under frozen generation controls
(temperature 0, top-p 1, seed 42, max 4096 tokens). We chose Qwen3.5-9B
because it is open-weight, carries first-class tool-calling support in vLLM
(native tool-call and reasoning parsers), and at 9B serves in bf16 on a
single 48\,GB card with 16k context. Using the same model for labeling and
serving keeps the label distribution matched to the engine under test. Three rows whose labeling
runs failed are declared as structural exclusions rather than silently
dropped, and three censored rows are excluded, leaving 5{,}994 eligible rows
split 3{,}997/998/999 by a per-row split column
fixed before any experiment; every run in this paper shares that split.
Label censoring is 0.05\%.

\subsection{Ranker and Feature Variants}
The Ranker is BERT-base~\cite{bert} producing a scalar score
$s_\theta(x) \in [0,1]$ from input rendered as
\texttt{[USER]}\,prompt\,\texttt{[TOOLS]}\,schema, truncated at 512 tokens.
We fine-tune with the pairwise margin objective
\begin{equation}
\mathcal{L}(x_a, x_b) = \max\bigl(0,\; -y \cdot (s_\theta(x_a) - s_\theta(x_b)) + m\bigr),
\end{equation}
where $y = \pm 1$ encodes which request produced the longer completion,
$m{=}1.0$, and pairs enter training only when their normalized length gap
exceeds $\delta{=}0.2$; three seeds (17/42/73) vary initialization and pair
shuffling over one fixed split. Table~\ref{tab:variants} lists the Feature Variants; the
full-context arm is reported void because the same 512-token cap leaves
80\% of its inputs truncated --- a disclosure, not a result.

\begin{table}[t]
  \caption{Feature variants and offline recipe. All variants share one fixed
  3{,}997/998/999 split of the labelled ToolACE sample.}
  \label{tab:variants}
  \centering\footnotesize
  \setlength{\tabcolsep}{3.5pt}
  \begin{tabular}{@{}lll@{}}
    \toprule
    Variant & Ranker input & Note \\
    \midrule
    prompt\_only & prompt text & encoder-strength control \\
    prompt\_schema & prompt + schema text & the proposed input \\
    full\_context & + history & void under truncation \\
    \midrule
    scalar (LightGBM) & 5 request statistics & fixed + 20-config grid \\
    identity (hash) & scalar + tool-set fingerprint & lookup emulation \\
    \bottomrule
  \end{tabular}
\end{table}

\subsection{Baselines and Identity Features}
The scalar baseline is LightGBM~\cite{lightgbm} over prompt length (chars,
words),
schema length, tool count, and history turns --- the scalar family
gateway-side length predictors describe --- in a fixed-hyperparameter form and a 20-configuration
grid search whose configuration is chosen on the validation split, never on
test; the stronger result is the baseline of record. The recipe is
deterministic with respect to seed (no stochastic sampling consumes the
random state), verified by bit-exact rerun. Identity baselines fingerprint
the tool set as SHA-256 over sorted tool names --- never over the raw schema
string, which embeds per-request timestamps that would fake uniqueness ---
and add out-of-fold per-fingerprint label statistics, emulating a deployed
per-client lookup table with a global-mean fallback for unseen keys.

\subsection{Cold-Start Strata and the Gate}
The fixed test split is partitioned into S1--S4 by fingerprint and
tool-name novelty against the training vocabulary. Strata under 100 rows
(S1$=$45, S2$=$78 at test) are withheld from headline claims; they may
inform design when disclosed with their size. The abstention set itself is
fixed on the validation split alone: S1 and S2 fall below the same 100-row
floor there ($n{=}74$ and $n{=}82$, against 472 and 370 for S3 and S4), so
the gate has no validation basis on which to vouch for them. Their test
sizes are reported for power, never used to set the rule. The Reliability Gate assigns each stratum the
lower bound of the 95\% session-clustered bootstrap~\cite{efron1979}
confidence interval of
\emph{validation} Ranking Tau --- fit on validation, evaluated on test ---
and abstains (confidence 0) on strata it cannot measure, including zero-tool
requests:
\begin{equation}
c(x) = \begin{cases}
  \max\bigl(0,\ \underline{\tau}_{95}(S(x))\bigr) & S(x) \in \{\mathrm{S3},\mathrm{S4}\} \\
  0 & \text{otherwise,}
\end{cases}
\end{equation}
where $S(x)$ is the request's stratum under the training vocabulary and
$\underline{\tau}_{95}$ the lower bootstrap bound of validation Ranking Tau.
A request is trusted when $c(x) \ge 0.5$; abstained requests keep their
arrival slots, and trusted requests reorder only among trusted slots. This confidence is a stratum-level abstention score, not a
calibrated probability.

\subsection{Serving Configuration}
The gateway runs in its default configuration plus exactly six
environment variables: listen address, API key, upstream provider and
model, the decision endpoint, and its timeout. Its request-shaping
pipeline, provider combos, semantic cache and scheduler-side queue
limits are all at their disabled defaults. The decision timeout is
6{,}000\,ms during ordering runs --- deliberately far above the
documented 15\,ms scheduling budget, so that fail-open never fires;
differences between ranked and unranked arms then bundle the ranking
signal with the synchronous scoring cost that carries it, and are read
as such. The budget itself is enforced and measured in dedicated arms
at 15\,ms (the documented default), 50 and 75\,ms
($\approx$1.3$\times$ and 2$\times$ the Ranker's GPU median), tracing
in-budget coverage against the operating point an operator who had
measured our scorer would actually set. vLLM prefix caching is disabled uniformly across the default-cap
ordering arms, verified from the engine's reported hit rate rather
than from the flag alone. It is enabled in the arm that measures its
effect and in all four arms of the capped-batch round, where holding it
on everywhere means it cannot separate them.

\subsection{Serving Environments and Provenance}
Table~\ref{tab:envs} pins both environments. Serving experiments
run on a single cloud RTX~4090 with 48\,GB VRAM (a vendor-modified memory
configuration common on Chinese GPU clouds); the decision service runs the
Ranker in fp16 with micro-batching (batch $\le$ 8, 3\,ms window) and a
reliability threshold of exactly 0.5. All vLLM capacity and cache settings
are held constant across every arm of every comparison; prefix caching is
disabled in every default-cap ordering arm --- verified from the engine's own
reported hit rate rather than from the flag alone --- and enabled in the
round that measures it and throughout the capped-batch round. The scheduler is
replaced through vLLM's~\cite{vllm2023} \texttt{--scheduler-cls} seam. Each
policy arm subclasses the engine's scheduler and re-sorts only the waiting
queue before each step, all else stock code; the stock pair differs in
scheduler implementation as well, the difference E5's attribution control
measures.

\begin{table}[t]
  \caption{Environment provenance for the two measurement platforms, which
  are separate rentals: labelling and fine-tuning ran first, serving weeks
  later.}
  \label{tab:envs}
  \centering\footnotesize
  \setlength{\tabcolsep}{3.5pt}
  \begin{tabular}{@{}lll@{}}
    \toprule
     & Training & Serving \\
    \midrule
    GPU & RTX 3090 24\,GB & RTX 4090 48\,GB (modified) \\
    Driver / CUDA & --- & 580.159.03 / 13.0 wheels \\
    Engine & vLLM 0.19.1, torch 2.10 & vLLM 0.24, torch 2.11 \\
    Model & Qwen3.5-9B (labels) & Qwen3.5-9B, bf16 \\
    Scheduler seam & --- & \texttt{--scheduler-cls} per arm \\
    \bottomrule
  \end{tabular}
\end{table}
